\documentclass[11pt,a4paper]{article}
\usepackage{fontspec,amsmath,amssymb,amsthm,microtype}
\setmainfont{lmroman10-regular.otf}[BoldFont=lmroman10-bold.otf,ItalicFont=lmroman10-italic.otf,BoldItalicFont=lmroman10-bolditalic.otf]
\usepackage{xcolor}
\usepackage[most]{tcolorbox}
\usepackage{mathrsfs}
\providecommand{\lt}{<}
\providecommand{\gt}{>}
\setmonofont{DejaVuSansMono.ttf}[BoldFont=DejaVuSansMono-Bold.ttf,ItalicFont=DejaVuSansMono-Oblique.ttf,BoldItalicFont=DejaVuSansMono-BoldOblique.ttf,Scale=0.86]
\newfontfamily\leanfallback{STIXTwoMath-Regular.otf}[Scale=0.86]
\definecolor{leanrule}{gray}{0.62}
\definecolor{leanback}{gray}{0.962}
\definecolor{leantitle}{gray}{0.22}
\definecolor{leanmuted}{gray}{0.38}
\newtcolorbox{leanbox}[1]{enhanced,breakable,colback=leanback,colframe=leanrule,
  boxrule=0.4pt,arc=0.6mm,left=2.4mm,right=2.4mm,top=1.8mm,bottom=1.8mm,
  before skip=0.8em,after skip=1em,fonttitle=\small\bfseries,coltitle=leantitle,
  colbacktitle=leanback,titlerule=0.3pt,title={#1}}
\newlength{\leancharwidth}
\newenvironment{leancode}{\par\smallskip\begingroup\ttfamily\footnotesize\raggedright
  \settowidth{\leancharwidth}{M}\setlength{\parskip}{0pt}\setlength{\parindent}{0pt}
  \frenchspacing}{\par\endgroup\smallskip}
\newcommand{\leanline}[2]{\hangindent=\dimexpr#1\leancharwidth+4\leancharwidth\relax
  \hangafter=1\noindent\hspace*{#1\leancharwidth}#2\par}
\newcommand{\leansource}[1]{\par{\footnotesize\color{leanmuted}#1}\par}
\newcommand{\leanlabel}[1]{\par\smallskip\noindent{\small\bfseries #1}\par}
\usepackage[margin=25mm]{geometry}
\usepackage[unicode,hidelinks]{hyperref}
\hypersetup{pdfauthor={LangevinLean Project}}
\setlength{\emergencystretch}{3em}
\setlength{\parskip}{0.4em}
\setlength{\parindent}{0pt}
\allowdisplaybreaks[2]
\raggedbottom
\title{\texorpdfstring{The entropy chain rule for the Gaussian interpolation}{The entropy chain rule for the Gaussian interpolation}}
\author{LangevinLean Project}
\date{}
\begin{document}\maketitle
Draft mathematical exposition; these notes may contain mistakes.


Let \(\pi=Z^{-1}e^{-U}dy\), let \(Y_t=X-t\nabla U(X)+\sqrt{2t}Z_0\) with independent standard Gaussian \(Z_0\), and write \(\rho_t=\mathcal L(Y_t)\), \(r_t=d\rho_t/d\pi\), and \(S_t=\nabla\log r_t\). Under the stated potential, log-Sobolev, and finite initial gradient-energy hypotheses, positive-time entropy is finite. Its time regularity is proved by integrating bounded approximations.



For the bounded tests \(\chi_m(\log r_t)\), the Gaussian weak equation and the time chain rule give exact integrated identities on \([a,b]\subset(0,\infty)\). The entropy integrand dominates their endpoints. Half the sum of the frozen-velocity and relative-score energies dominates their fluxes independently of \(m\) and is integrable in time. Dominated convergence therefore yields \(H(\rho_b\mid\pi)-H(\rho_a\mid\pi)=\int_a^bq(t)dt\). This proves local absolute continuity and derivative \(q\) almost everywhere. The score-pairing identity subsequently writes that derivative as \(-I_t+\mathbb E\langle\nabla U(Y_t)-\nabla U(X),S_t(Y_t)\rangle\).



Earlier results used here: \cite[Section~3.46]{LangevinLean2026ULASC}.

\section{\texorpdfstring{Removing bounded log weights by dominated convergence}{Removing bounded log weights by dominated convergence}}


Let \(w_m(z)=[1+(z/(m+1))^2]^{-1}\). For a measurable scalar field \(z\) and an integrable field \(f\),

\[\int w_m(z(x))f(x)\,d\nu(x)\longrightarrow\int f(x)\,d\nu(x).\]

The pointwise limit is \(w_m(z(x))\to1\), proved in the pointwise convergence of bounded-log derivative weights to one. Each weighted integrand is measurable, because the weight is obtained by division, squaring and inversion of measurable real functions. Furthermore \(0\le w_m\le1\), so \(|w_m(z)f|\le|f|\). The proof invokes dominated convergence with the integrable majorant \(|f|\). This lemma later removes the truncation from the Gaussian flux, whose two factors are square integrable.



The derivative weights converge pointwise to one and stay between zero and one. Multiplying by the already integrable unweighted pairing gives a common absolute majorant, so dominated convergence removes the weight without changing the reference measure.

\begin{leanbox}{Lean: \texttt{boundedLogWeight\_\allowbreak{}integral\_\allowbreak{}tendsto}}
\leansource{\texttt{boundedLogWeight\_\allowbreak{}integral\_\allowbreak{}tendsto} \quad UnadjustedLangevin/\allowbreak{}LSI/\allowbreak{}GaussianEntropyIdentity.lean, lines 30–46}
\begin{leancode}
\leanline{0}{theorem\ boundedLogWeight\_integral\_tendsto}
\leanline{4}{(ν\ :\ Measure\ Ω)\ (z\ f\ :\ Ω\ →\ ℝ)\ (hz\ :\ Measurable\ z)\ (hf\ :\ Integrable\ f\ ν)\ :}
\leanline{4}{Tendsto\ (fun\ m\ :\ ℕ\ ↦\ ∫\ x,\ boundedLogWeight\ m\ (z\ x)\ *\ f\ x\ ∂ν)}
\leanline{6}{atTop\ ({\leanfallback 𝓝}\ (∫\ x,\ f\ x\ ∂ν))}
\end{leancode}
\leanlabel{How the Lean proof follows the argument}
\texttt{bounded\allowbreak{}Log\allowbreak{}Weight\_\allowbreak{}integral\_\allowbreak{}tendsto}: the intermediate fact is \texttt{hm}, in source order; the source steps apply or rewrite \texttt{tendsto\_\allowbreak{}integral\_\allowbreak{}of\_\allowbreak{}dominated\_\allowbreak{}convergence}, \texttt{hm.\allowbreak{}aestrongly\allowbreak{}Measurable.\allowbreak{}mul}, \texttt{one\_\allowbreak{}mul}, \texttt{bounded\allowbreak{}Log\allowbreak{}Weight\_\allowbreak{}tendsto}, \texttt{mul\_\allowbreak{}const}.
\end{leanbox}

\section{\texorpdfstring{Why the square-integrable flux pairing is integrable}{Why the square-integrable flux pairing is integrable}}


If \(v,s\in L^2(\nu;E)\), then

\[|\langle v,s\rangle|\le\|v\|\|s\|,\qquad \|v\|\|s\|\in L^1(\nu).\]

Cauchy–Schwarz bounds the integral of the product of the two norms by their finite square norms. The absolute inner product is bounded by this norm product, proving absolute integrability of the pairing. This remains true on the stated measure space without assuming its total mass is finite.



For square-integrable fields \(a,b\), use \(|\langle a,b\rangle|\le(\|a\|^2+\|b\|^2)/2\). The right side has finite integral, proving absolute flux integrability before signed Fubini or limiting arguments.

\begin{leanbox}{Lean: \texttt{integrable\_\allowbreak{}inner\_\allowbreak{}of\_\allowbreak{}memLp\_\allowbreak{}two}}
\leansource{\texttt{integrable\_\allowbreak{}inner\_\allowbreak{}of\_\allowbreak{}memLp\_\allowbreak{}two} \quad UnadjustedLangevin/\allowbreak{}LSI/\allowbreak{}GaussianEntropyIdentity.lean, lines 52–59}
\begin{leancode}
\leanline{0}{theorem\ integrable\_inner\_of\_memLp\_two\ (ν\ :\ Measure\ Ω)\ (f\ g\ :\ Ω\ →\ E)}
\leanline{4}{(hf\ :\ MemLp\ f\ 2\ ν)\ (hg\ :\ MemLp\ g\ 2\ ν)\ :}
\leanline{4}{Integrable\ (fun\ x\ ↦\ inner\ ℝ\ (f\ x)\ (g\ x))\ ν}
\end{leancode}
\leanlabel{How the Lean proof follows the argument}
\texttt{integrable\_\allowbreak{}inner\_\allowbreak{}of\_\allowbreak{}mem\allowbreak{}Lp\_\allowbreak{}two}: the intermediate fact is \texttt{hp}, in source order; the source steps apply or rewrite \texttt{hp.\allowbreak{}mono'}, \texttt{hf.\allowbreak{}aestrongly\allowbreak{}Measurable.\allowbreak{}inner}, \texttt{hg.\allowbreak{}aestrongly\allowbreak{}Measurable}, \texttt{ae\_\allowbreak{}of\_\allowbreak{}all}, \texttt{norm\_\allowbreak{}inner\_\allowbreak{}le\_\allowbreak{}norm}.
\end{leanbox}

\section{\texorpdfstring{The real Gibbs density and its normalization}{The real Gibbs density and its normalization}}


Write \(Z=\int e^{-U(x)}\,dx\), and let \(\pi\) be the normalized Gibbs law. Define

\[J(y)=Z^{-1}e^{-U(y)}.\]

The existing probability-mass argument proves \(0<Z<\infty\); differentiability makes \(U\) measurable and gives strict positivity of the partition function. Consequently \(J(y)>0\) at every point. The measure identity is

\[\pi(dy)=J(y)\,dy.\]

The density identity follows from multiplying the unnormalized Gibbs weight by the finite scalar \(Z^{-1}\). Since \(0<Z<\infty\), the resulting measurable function \(J\) is positive and has integral one. The same expression therefore represents the target in ordinary Lebesgue integration; strong convexity is not needed for this change-of-density step.



The Gibbs density is positive and finite pointwise, and its integral is one. Thus \(J\Phi_m(r_t)\) may be read as a Lebesgue representation of the same bounded relative-entropy primitive integrated against Gibbs.

\begin{leanbox}{Lean: \texttt{gibbsLebesgueDensity}, \texttt{gibbsLebesgueDensity\_\allowbreak{}pos}, \texttt{gibbsMeasure\_\allowbreak{}eq\_\allowbreak{}withDensity\_\allowbreak{}real}}
\leansource{\texttt{gibbsLebesgueDensity} \quad UnadjustedLangevin/\allowbreak{}LSI/\allowbreak{}GaussianEntropyIdentity.lean, lines 63–64}
\begin{leancode}
\leanline{0}{def\ gibbsLebesgueDensity\ (U\ :\ E\ →\ ℝ)\ (y\ :\ E)\ :\ ℝ\ :=}
\leanline{2}{(ULA.gibbsPartition\ U).toReal⁻¹\ *\ Real.exp\ (-U\ y)}
\end{leancode}
\leansource{\texttt{gibbsLebesgueDensity\_\allowbreak{}pos} \quad UnadjustedLangevin/\allowbreak{}LSI/\allowbreak{}GaussianEntropyIdentity.lean, lines 66–70}
\begin{leancode}
\leanline{0}{theorem\ gibbsLebesgueDensity\_pos\ (U\ :\ E\ →\ ℝ)\ (hU\ :\ Differentiable\ ℝ\ U)}
\leanline{4}{[IsProbabilityMeasure\ (ULA.gibbsMeasure\ U)]\ (y\ :\ E)\ :}
\leanline{4}{0\ <\ gibbsLebesgueDensity\ U\ y}
\end{leancode}
\leansource{\texttt{gibbsMeasure\_\allowbreak{}eq\_\allowbreak{}withDensity\_\allowbreak{}real} \quad UnadjustedLangevin/\allowbreak{}LSI/\allowbreak{}GaussianEntropyIdentity.lean, lines 74–87}
\begin{leancode}
\leanline{0}{theorem\ gibbsMeasure\_eq\_withDensity\_real\ (U\ :\ E\ →\ ℝ)\ (hU\ :\ Differentiable\ ℝ\ U)}
\leanline{4}{[IsProbabilityMeasure\ (ULA.gibbsMeasure\ U)]\ :}
\leanline{4}{ULA.gibbsMeasure\ U\ =\ volume.withDensity\ (fun\ y\ ↦}
\leanline{6}{ENNReal.ofReal\ (gibbsLebesgueDensity\ U\ y))}
\end{leancode}
\leanlabel{How the Lean proof follows the argument}
\texttt{gibbs\allowbreak{}Lebesgue\allowbreak{}Density} is a definition with defining expression \texttt{(ULA.\allowbreak{}gibbs\allowbreak{}Partition U).\allowbreak{}to\allowbreak{}Real\ensuremath{{}^{-1}} * Real.\allowbreak{}exp (-\allowbreak{}U y)}. The right-hand side specifies the object used in the mathematical statement.



\texttt{gibbs\allowbreak{}Lebesgue\allowbreak{}Density\_\allowbreak{}pos}: the source steps apply or rewrite \texttt{mul\_\allowbreak{}pos}, \texttt{inv\_\allowbreak{}pos.\allowbreak{}mpr}, \texttt{ENNReal.\allowbreak{}to\allowbreak{}Real\_\allowbreak{}pos}, \texttt{ULA.\allowbreak{}gibbs\allowbreak{}Partition\_\allowbreak{}pos}.



\texttt{gibbs\allowbreak{}Measure\_\allowbreak{}eq\_\allowbreak{}with\allowbreak{}Density\_\allowbreak{}real}: the intermediate fact is \texttt{hw}, in source order; the source steps apply or rewrite \texttt{ULA.\allowbreak{}gibbs\allowbreak{}Measure}, \texttt{with\allowbreak{}Density\_\allowbreak{}smul}, \texttt{ENNReal.\allowbreak{}of\allowbreak{}Real\_\allowbreak{}mul}, \texttt{ENNReal.\allowbreak{}to\allowbreak{}Real\_\allowbreak{}inv}, \texttt{ENNReal.\allowbreak{}of\allowbreak{}Real\_\allowbreak{}to\allowbreak{}Real}.
\end{leanbox}


Earlier results used here: \cite[Section~3.26.1]{LangevinLean2026ULASC}; \cite[Section~3.26.2]{LangevinLean2026ULASC}; \cite[Section~3.26.3]{LangevinLean2026ULASC}; \cite[Section~3.26.9]{LangevinLean2026ULASC}.

\section{\texorpdfstring{The time derivative of the weighted entropy primitive}{The time derivative of the weighted entropy primitive}}


Let \(p_t\) be the Lebesgue density of the frozen Gaussian update, and let \(r_t=p_t/J\) be its relative density. For the bounded entropy primitive \(\Phi_m\), one has \(\Phi_m'(r)=\chi_m(\log r)\) for \(r>0\). At every positive time, the kernel differentiation result gives \(\partial_t p_t=\dot p_t\), hence

\[\partial_t r_t=Ze^{U}\dot p_t,\qquad \partial_t[J\Phi_m(r_t)]=\chi_m(\log r_t)\dot p_t.\]

The proof applies the time chain rule first to the constant multiple defining \(r_t\), then to the derivative formula for the bounded entropy primitive, then to multiplication by \(J\). The remaining coefficient cancels because \(J(y)Ze^{U(y)}=1\), using the inverse of the strictly positive real partition function and \(e^{-U}e^U=1\). This cancellation is what permits the spacetime estimate to use the Lebesgue \(L^1\) bound on \(\dot p_t\).



Because \(J\) is time independent and \(Jr_t=p_t\), differentiating \(J\Phi_m(r_t)\) gives \(\chi_m(\log r_t)\dot p_t\). This cancellation avoids differentiating the Gibbs density in time.

\begin{leanbox}{Lean: \texttt{hasDerivAt\_\allowbreak{}weighted\_\allowbreak{}boundedEntropy}}
\leansource{\texttt{hasDerivAt\_\allowbreak{}weighted\_\allowbreak{}boundedEntropy} \quad UnadjustedLangevin/\allowbreak{}LSI/\allowbreak{}GaussianEntropyIdentity.lean, lines 92–127}
\begin{leancode}
\leanline{0}{theorem\ hasDerivAt\_weighted\_boundedEntropy}
\leanline{4}{(U\ :\ E\ →\ ℝ)\ (μ\ :\ Measure\ E)\ [IsProbabilityMeasure\ μ]}
\leanline{4}{(L\ :\ ℝ≥0)\ (hU\ :\ Differentiable\ ℝ\ U)\ (hLip\ :\ LipschitzWith\ L\ (gradient\ U))}
\leanline{4}{[IsProbabilityMeasure\ (ULA.gibbsMeasure\ U)]}
\leanline{4}{(hG\ :\ Integrable\ (fun\ x\ ↦\ ‖gradient\ U\ x‖\ \textasciicircum{}\ 2)\ μ)}
\leanline{4}{(t\ :\ ℝ)\ (ht\ :\ 0\ <\ t)\ (m\ :\ ℕ)\ (y\ :\ E)\ :}
\leanline{4}{HasDerivAt\ (fun\ s\ ↦\ gibbsLebesgueDensity\ U\ y\ *}
\leanline{6}{boundedEntropy\ m\ (interpolationRelativeDensity\ U\ μ\ s\ y))}
\leanline{6}{(boundedLogSlope\ m\ (Real.log\ (interpolationRelativeDensity\ U\ μ\ t\ y))\ *}
\leanline{8}{interpolationLebesgueDensityTimeDerivative\ U\ μ\ t\ y)\ t}
\end{leancode}
\leanlabel{How the Lean proof follows the argument}
\texttt{has\allowbreak{}Deriv\allowbreak{}At\_\allowbreak{}weighted\_\allowbreak{}bounded\allowbreak{}Entropy}: the intermediate facts include \texttt{htime}, \texttt{hrel}, \texttt{hcomp}, \texttt{h\allowbreak{}Z}, \texttt{hcancel}, in source order; the source steps apply or rewrite \texttt{htime.\allowbreak{}const\_\allowbreak{}mul}, \texttt{inv\_\allowbreak{}mul\_\allowbreak{}cancel\ensuremath{{}_{0}}}, \texttt{Real.\allowbreak{}exp\_\allowbreak{}zero}, \texttt{mul\_\allowbreak{}one}, \texttt{one\_\allowbreak{}mul}. The \texttt{calc} block checks the successive equalities or inequalities in the argument. The remaining scalar rearrangements are checked by \texttt{ring}.
\end{leanbox}


Earlier results used here: \cite[Section~3.26.2]{LangevinLean2026ULASC}; \cite[Section~3.26.3]{LangevinLean2026ULASC}; \cite[Section~3.26.9]{LangevinLean2026ULASC}.

\section{\texorpdfstring{Identifying finite entropy through the chosen density}{Identifying finite entropy through the chosen density}}


Let \(\nu\) be the reference probability and suppose \(\mu=r\nu\) is a probability measure, \(r\ge0\) is measurable, and \(H(\mu\mid\nu)<\infty\). Then

\[\mathrm{kl}(r)=r\log r+1-r\in L^1(\nu),\qquad \int\mathrm{kl}(r)\,d\nu=H(\mu\mid\nu).\]

The Radon–Nikodym derivative of \(r\nu\) relative to \(\nu\) is \(r\) almost everywhere. Finite KL gives integrability of its logarithm under \(r\nu\); changing measure makes \(r\log r\) integrable under \(\nu\). Since \(r\) and the constant one both have finite integral, so does \(r\log r-r+1\). Their probability masses agree, so integrating the last two terms cancels them and gives the stated entropy identity.



The relative entropy is finite at positive time by LSI and the Fisher bound. The density formula therefore identifies it with the ordinary integral of \(r_t\log r_t-r_t+1\), whose mass corrections cancel by normalization.

\begin{leanbox}{Lean: \texttt{integrable\_\allowbreak{}klFun\_\allowbreak{}density\_\allowbreak{}of\_\allowbreak{}finite\_\allowbreak{}kl}}
\leansource{\texttt{integrable\_\allowbreak{}klFun\_\allowbreak{}density\_\allowbreak{}of\_\allowbreak{}finite\_\allowbreak{}kl} \quad UnadjustedLangevin/\allowbreak{}LSI/\allowbreak{}GaussianEntropyIdentity.lean, lines 131–149}
\begin{leancode}
\leanline{0}{theorem\ integrable\_klFun\_density\_of\_finite\_kl}
\leanline{4}{(ν\ :\ Measure\ E)\ [IsProbabilityMeasure\ ν]\ (r\ :\ E\ →\ ℝ)}
\leanline{4}{(hr\ :\ Measurable\ r)\ (hpos\ :\ ∀\ y,\ 0\ ≤\ r\ y)}
\leanline{4}{[IsProbabilityMeasure\ (tiltedMeasure\ ν\ r)]}
\leanline{4}{(hKL\ :\ InformationTheory.klDiv\ (tiltedMeasure\ ν\ r)\ ν\ ≠\ ⊤)\ :}
\leanline{4}{Integrable\ (fun\ y\ ↦\ InformationTheory.klFun\ (r\ y))\ ν\ ∧}
\leanline{6}{(∫\ y,\ InformationTheory.klFun\ (r\ y)\ ∂ν)\ =}
\leanline{8}{(InformationTheory.klDiv\ (tiltedMeasure\ ν\ r)\ ν).toReal}
\end{leancode}
\leanlabel{How the Lean proof follows the argument}
\texttt{integrable\_\allowbreak{}kl\allowbreak{}Fun\_\allowbreak{}density\_\allowbreak{}of\_\allowbreak{}finite\_\allowbreak{}kl}: the intermediate facts include \texttt{hac}, \texttt{hrn}, \texttt{heq}, \texttt{hi}, in source order; the source steps apply or rewrite \texttt{ENNReal.\allowbreak{}to\allowbreak{}Real\_\allowbreak{}of\allowbreak{}Real}, \texttt{hi.\allowbreak{}congr}, \texttt{Information\allowbreak{}Theory.\allowbreak{}to\allowbreak{}Real\_\allowbreak{}kl\allowbreak{}Div\_\allowbreak{}eq\_\allowbreak{}integral\_\allowbreak{}kl\allowbreak{}Fun}, \texttt{integral\_\allowbreak{}congr\_\allowbreak{}ae}.
\end{leanbox}

\section{\texorpdfstring{Integrability of every bounded entropy approximation}{Integrability of every bounded entropy approximation}}


At \(t>0\), the interpolation has a positive \(C^1\) relative density \(r_t\). Its Fisher information is finite from the initial score moment and Gaussian smoothing; LSI then gives finite relative entropy. The bounded primitive satisfies

\[0\le\Phi_m(r_t(y))\le\mathrm{kl}(r_t(y)).\]

The preceding density lemma proves that the right-hand side is integrable. The left-hand side is measurable: continuity of \(r_t\) and the derivative of \(\Phi_m\) at positive arguments give continuity of their composition. Integrable domination therefore proves integrability for each truncation from this finite entropy majorant. Initial density differentiability, initial Fisher information and a second moment of the state are not required.



The bounded primitive lies between zero and the nonnegative entropy integrand. Finite relative entropy therefore proves integrability of each approximating entropy before its time identity is formed.

\begin{leanbox}{Lean: \texttt{interpolation\_\allowbreak{}boundedEntropy\_\allowbreak{}integrable}}
\leansource{\texttt{interpolation\_\allowbreak{}boundedEntropy\_\allowbreak{}integrable} \quad UnadjustedLangevin/\allowbreak{}LSI/\allowbreak{}GaussianEntropyIdentity.lean, lines 153–182}
\begin{leancode}
\leanline{0}{theorem\ interpolation\_boundedEntropy\_integrable}
\leanline{4}{(U\ :\ E\ →\ ℝ)\ (μ\ :\ Measure\ E)\ [IsProbabilityMeasure\ μ]}
\leanline{4}{(α\ :\ ℝ)\ (hα\ :\ 0\ <\ α)\ (L\ :\ ℝ≥0)\ (hU\ :\ Differentiable\ ℝ\ U)}
\leanline{4}{(hLip\ :\ LipschitzWith\ L\ (gradient\ U))}
\leanline{4}{(hLSI\ :\ IsLogSobolev\ (ULA.gibbsMeasure\ U)\ α)}
\leanline{4}{(hgrad\ :\ MemLp\ (gradient\ U)\ 2\ μ)\ (t\ :\ ℝ)\ (ht\ :\ 0\ <\ t)\ (m\ :\ ℕ)\ :}
\leanline{4}{Integrable\ (fun\ y\ ↦\ boundedEntropy\ m\ (interpolationRelativeDensity\ U\ μ\ t\ y))}
\leanline{6}{(ULA.gibbsMeasure\ U)}
\end{leancode}
\leanlabel{How the Lean proof follows the argument}
\texttt{interpolation\_\allowbreak{}bounded\allowbreak{}Entropy\_\allowbreak{}integrable}: the intermediate facts include \texttt{hpos}, \texttt{h\allowbreak{}KL}, \texttt{hi}, \texttt{hcont}, in source order; the source steps apply or rewrite \texttt{tilted\allowbreak{}Measure\_\allowbreak{}interpolation\allowbreak{}Relative\allowbreak{}Density}, \texttt{interpolation\allowbreak{}Law\_\allowbreak{}kl\allowbreak{}Div\_\allowbreak{}ne\_\allowbreak{}top}, \texttt{abs\_\allowbreak{}of\_\allowbreak{}nonneg}, \texttt{bounded\allowbreak{}Entropy\_\allowbreak{}nonneg}, \texttt{bounded\allowbreak{}Entropy\_\allowbreak{}le\_\allowbreak{}kl\allowbreak{}Fun}.
\end{leanbox}


Earlier results used here: \cite[Section~3.26.3]{LangevinLean2026ULASC}.

\section{\texorpdfstring{Convergence of the bounded entropy to relative entropy}{Convergence of the bounded entropy to relative entropy}}


For every fixed positive time,

\[\int\Phi_m(r_t)\,d\pi\longrightarrow H(\rho_t\mid\pi).\]

Pointwise convergence \(\Phi_m(r)\to\mathrm{kl}(r)\) holds for \(r>0\). The domination \(0\le\Phi_m(r_t)\le\mathrm{kl}(r_t)\), together with integrability proved above, gives dominated convergence under the fixed Gibbs measure. The density identity \(r_t\pi=\rho_t\), with \(r_t=Ze^Up_t\), identifies the integral of \(r_t\log r_t-r_t+1\) with \(H(\rho_t\mid\pi)\). The limit is therefore the relative entropy of the specified Gaussian interpolation law.



Pointwise convergence of the primitive to \(r\log r-r+1\), together with the finite entropy majorant, gives convergence of both positive-time endpoint entropies. This limit is independent of any time-flux limit.

\begin{leanbox}{Lean: \texttt{interpolation\_\allowbreak{}boundedEntropy\_\allowbreak{}tendsto}}
\leansource{\texttt{interpolation\_\allowbreak{}boundedEntropy\_\allowbreak{}tendsto} \quad UnadjustedLangevin/\allowbreak{}LSI/\allowbreak{}GaussianEntropyIdentity.lean, lines 187–222}
\begin{leancode}
\leanline{0}{theorem\ interpolation\_boundedEntropy\_tendsto}
\leanline{4}{(U\ :\ E\ →\ ℝ)\ (μ\ :\ Measure\ E)\ [IsProbabilityMeasure\ μ]}
\leanline{4}{(α\ :\ ℝ)\ (hα\ :\ 0\ <\ α)\ (L\ :\ ℝ≥0)\ (hU\ :\ Differentiable\ ℝ\ U)}
\leanline{4}{(hLip\ :\ LipschitzWith\ L\ (gradient\ U))}
\leanline{4}{(hLSI\ :\ IsLogSobolev\ (ULA.gibbsMeasure\ U)\ α)}
\leanline{4}{(hgrad\ :\ MemLp\ (gradient\ U)\ 2\ μ)\ (t\ :\ ℝ)\ (ht\ :\ 0\ <\ t)\ :}
\leanline{4}{Tendsto\ (fun\ m\ :\ ℕ\ ↦\ ∫\ y,}
\leanline{6}{boundedEntropy\ m\ (interpolationRelativeDensity\ U\ μ\ t\ y)\ ∂ULA.gibbsMeasure\ U)}
\leanline{6}{atTop\ ({\leanfallback 𝓝}\ ((InformationTheory.klDiv\ (interpolationLaw\ U\ μ\ t)}
\leanline{8}{(ULA.gibbsMeasure\ U)).toReal))}
\end{leancode}
\leanlabel{How the Lean proof follows the argument}
\texttt{interpolation\_\allowbreak{}bounded\allowbreak{}Entropy\_\allowbreak{}tendsto}: the intermediate facts include \texttt{hpos}, \texttt{h\allowbreak{}KL}, in source order; the source steps apply or rewrite \texttt{tilted\allowbreak{}Measure\_\allowbreak{}interpolation\allowbreak{}Relative\allowbreak{}Density}, \texttt{interpolation\allowbreak{}Law\_\allowbreak{}kl\allowbreak{}Div\_\allowbreak{}ne\_\allowbreak{}top}, \texttt{bounded\allowbreak{}Entropy\_\allowbreak{}nonneg}, \texttt{bounded\allowbreak{}Entropy\_\allowbreak{}le\_\allowbreak{}kl\allowbreak{}Fun}, \texttt{bounded\allowbreak{}Entropy\_\allowbreak{}tendsto}.
\end{leanbox}


Earlier results used here: \cite[Section~3.26.3]{LangevinLean2026ULASC}.

\section{\texorpdfstring{The bounded entropy and its time flux}{The bounded entropy and its time flux}}


The two functions used in the time integration are

\[H_m(t)=\int\Phi_m(r_t(y))\,d\pi(y),\qquad q_m(t)=\int\chi_m(\log r_t(y))\dot p_t(y)\,dy.\]

The bounded entropy \(H_m(t)\) uses the fixed Gibbs reference measure; the bounded entropy flux \(q_m(t)\) uses Lebesgue measure. Their different measures match the exact cancellation in the weighted chain rule. The definitions are explicit functions of the mixture density and its kernel time derivative.



The bounded entropy is paired with its explicit time flux on a compact positive interval. Boundedness of the log test lets the spacetime \(L^1\) density-derivative estimate dominate that flux.

\begin{leanbox}{Lean: \texttt{interpolationBoundedEntropy}, \texttt{interpolationBoundedEntropyFlux}}
\leansource{\texttt{interpolationBoundedEntropy} \quad UnadjustedLangevin/\allowbreak{}LSI/\allowbreak{}GaussianEntropyIdentity.lean, lines 225–226}
\begin{leancode}
\leanline{0}{def\ interpolationBoundedEntropy\ (U\ :\ E\ →\ ℝ)\ (μ\ :\ Measure\ E)\ (m\ :\ ℕ)\ (t\ :\ ℝ)\ :\ ℝ\ :=}
\leanline{2}{∫\ y,\ boundedEntropy\ m\ (interpolationRelativeDensity\ U\ μ\ t\ y)\ ∂ULA.gibbsMeasure\ U}
\end{leancode}
\leansource{\texttt{interpolationBoundedEntropyFlux} \quad UnadjustedLangevin/\allowbreak{}LSI/\allowbreak{}GaussianEntropyIdentity.lean, lines 229–231}
\begin{leancode}
\leanline{0}{def\ interpolationBoundedEntropyFlux\ (U\ :\ E\ →\ ℝ)\ (μ\ :\ Measure\ E)\ (m\ :\ ℕ)\ (t\ :\ ℝ)\ :\ ℝ\ :=}
\leanline{2}{∫\ y,\ boundedLogSlope\ m\ (Real.log\ (interpolationRelativeDensity\ U\ μ\ t\ y))\ *}
\leanline{4}{interpolationLebesgueDensityTimeDerivative\ U\ μ\ t\ y}
\end{leancode}
\leanlabel{How the Lean proof follows the argument}
\texttt{interpolation\allowbreak{}Bounded\allowbreak{}Entropy} is a definition with defining expression \texttt{\ensuremath{\int} y, bounded\allowbreak{}Entropy m (interpolation\allowbreak{}Relative\allowbreak{}Density U \ensuremath{\mu} t y) \ensuremath{\partial}ULA.\allowbreak{}gibbs\allowbreak{}Measure U}. The right-hand side specifies the object used in the mathematical statement.



\texttt{interpolation\allowbreak{}Bounded\allowbreak{}Entropy\allowbreak{}Flux} is a definition with defining expression \texttt{\ensuremath{\int} y, bounded\allowbreak{}Log\allowbreak{}Slope m (Real.\allowbreak{}log (interpolation\allowbreak{}Relative\allowbreak{}Density U \ensuremath{\mu} t y)) *}. The right-hand side specifies the object used in the mathematical statement.
\end{leanbox}


Earlier results used here: \cite[Section~3.26.3]{LangevinLean2026ULASC}.

\section{\texorpdfstring{Admissibility of the bounded logarithmic test}{Admissibility of the bounded logarithmic test}}


The spatial weak test is \(\psi_m(y)=\chi_m(\log r_t(y))\). Positivity and \(C^1\) regularity of \(r_t\), combined with smoothness of \(\chi_m\), show that \(\psi_m\in C^1\). Its uniform bound is

\[|\psi_m(y)|\le(m+1)\pi/2.\]

The spatial chain rule gives

\[\nabla\psi_m(y)=w_m(\log r_t(y))\nabla\log r_t(y),\qquad \|\nabla\psi_m(y)\|\le\|\nabla\log r_t(y)\|.\]

The test-gradient norm is at most the relative-score norm, whose square is integrable. Continuity of the test derivative supplies measurability. Thus the bounded logarithmic test is \(C^1\) and has finite gradient energy, exactly the conditions required by the Gaussian weak equation.



The test is bounded and \(C^1\), and its gradient norm is at most the relative-score norm. Finite Fisher information therefore places its gradient in the exact finite-energy class admitted by the weak equation.

\begin{leanbox}{Lean: \texttt{interpolationBoundedLogTest\_\allowbreak{}admissible}}
\leansource{\texttt{interpolationBoundedLogTest\_\allowbreak{}admissible} \quad UnadjustedLangevin/\allowbreak{}LSI/\allowbreak{}GaussianEntropyIdentity.lean, lines 235–258}
\begin{leancode}
\leanline{0}{theorem\ interpolationBoundedLogTest\_admissible}
\leanline{4}{(U\ :\ E\ →\ ℝ)\ (μ\ :\ Measure\ E)\ [IsProbabilityMeasure\ μ]}
\leanline{4}{(L\ :\ ℝ≥0)\ (hU\ :\ Differentiable\ ℝ\ U)\ (hLip\ :\ LipschitzWith\ L\ (gradient\ U))}
\leanline{4}{[IsProbabilityMeasure\ (ULA.gibbsMeasure\ U)]}
\leanline{4}{(hgrad\ :\ MemLp\ (gradient\ U)\ 2\ μ)\ (t\ :\ ℝ)\ (ht\ :\ 0\ <\ t)\ (m\ :\ ℕ)\ :}
\leanline{4}{ContDiff\ ℝ\ 1\ (fun\ y\ ↦\ boundedLogSlope\ m\ (Real.log\ (interpolationRelativeDensity\ U\ μ\ t\ y)))\ ∧}
\leanline{6}{(∀\ y,\ ‖boundedLogSlope\ m\ (Real.log\ (interpolationRelativeDensity\ U\ μ\ t\ y))‖\ ≤}
\leanline{8}{(m\ +\ 1\ :\ ℝ)\ *\ (Real.pi\ /\ 2))\ ∧}
\leanline{6}{MemLp\ (gradient\ (fun\ y\ ↦\ boundedLogSlope\ m}
\leanline{8}{(Real.log\ (interpolationRelativeDensity\ U\ μ\ t\ y))))\ 2\ (interpolationLaw\ U\ μ\ t)}
\end{leancode}
\leanlabel{How the Lean proof follows the argument}
\texttt{interpolation\allowbreak{}Bounded\allowbreak{}Log\allowbreak{}Test\_\allowbreak{}admissible}: the intermediate facts include \texttt{hpos}, \texttt{hlog}, \texttt{h\ensuremath{\psi}}, \texttt{hm}, in source order; the source steps apply or rewrite \texttt{bounded\allowbreak{}Log\allowbreak{}Slope\_\allowbreak{}cont\allowbreak{}Diff}, \texttt{of\_\allowbreak{}le}, \texttt{interpolation\allowbreak{}Relative\allowbreak{}Score\_\allowbreak{}mem\allowbreak{}Lp}, \texttt{ae\_\allowbreak{}of\_\allowbreak{}all}, \texttt{norm\_\allowbreak{}gradient\_\allowbreak{}bounded\allowbreak{}Log\allowbreak{}Slope\_\allowbreak{}comp\_\allowbreak{}le}.
\end{leanbox}


Earlier results used here: \cite[Section~3.26.3]{LangevinLean2026ULASC}.

\section{\texorpdfstring{The bounded entropy flux equals the spatial pairing}{The bounded entropy flux equals the spatial pairing}}


Let \(K_t(x,y)\) be the Gaussian transition density, \(g=\nabla U\), and \(S_t=\nabla\log r_t\). The proved Gaussian weak equation, applied to the admissible test above, gives

\[q_m(t)=-\iint K_t(x,y)w_m(\log r_t(y))\left\langle g(x)-\frac{y-x+tg(x)}{2t},S_t(y)\right\rangle\,\mu(dx)\,dy.\]

The underlying flux is \(K_t[g(x)-(y-x+tg(x))/(2t)]\). Its divergence equals the kernel time derivative; global integration by parts is justified by the previously proved finite-gradient-energy estimates. Substitute the test-gradient chain rule and the equality between the explicit measurable score field and the relative score. The remaining rearrangement is bilinearity of the inner product and scalar multiplication.



The weak equation converts the time-derivative pairing into the negative spatial flux pairing. Its sign is fixed by integration by parts, and its integrability has been established for this very bounded test.

\begin{leanbox}{Lean: \texttt{interpolationBoundedEntropyFlux\_\allowbreak{}eq\_\allowbreak{}joint}}
\leansource{\texttt{interpolationBoundedEntropyFlux\_\allowbreak{}eq\_\allowbreak{}joint} \quad UnadjustedLangevin/\allowbreak{}LSI/\allowbreak{}GaussianEntropyIdentity.lean, lines 263–296}
\begin{leancode}
\leanline{0}{theorem\ interpolationBoundedEntropyFlux\_eq\_joint}
\leanline{4}{(U\ :\ E\ →\ ℝ)\ (μ\ :\ Measure\ E)\ [IsProbabilityMeasure\ μ]}
\leanline{4}{(L\ :\ ℝ≥0)\ (hU\ :\ Differentiable\ ℝ\ U)\ (hLip\ :\ LipschitzWith\ L\ (gradient\ U))}
\leanline{4}{[IsProbabilityMeasure\ (ULA.gibbsMeasure\ U)]}
\leanline{4}{(hgrad\ :\ MemLp\ (gradient\ U)\ 2\ μ)\ (t\ :\ ℝ)\ (ht\ :\ 0\ <\ t)\ (m\ :\ ℕ)\ :}
\leanline{4}{interpolationBoundedEntropyFlux\ U\ μ\ m\ t\ =}
\leanline{6}{-(∫\ z\ :\ E\ ×\ E,}
\leanline{8}{realGaussianKernel\ (1\ /\ (4\ *\ t))\ (z.1\ -\ t\ •\ gradient\ U\ z.1)\ z.2\ *}
\leanline{10}{boundedLogWeight\ m\ (Real.log\ (interpolationRelativeDensity\ U\ μ\ t\ z.2))\ *}
\leanline{10}{inner\ ℝ\ (gradient\ U\ z.1\ -\ (1\ /\ (2\ *\ t))\ •}
\leanline{12}{(z.2\ -\ (z.1\ -\ t\ •\ gradient\ U\ z.1)))}
\leanline{12}{(interpolationRelativeScoreField\ U\ μ\ t\ z.2)\ ∂μ.prod\ volume)}
\end{leancode}
\leanlabel{How the Lean proof follows the argument}
\texttt{interpolation\allowbreak{}Bounded\allowbreak{}Entropy\allowbreak{}Flux\_\allowbreak{}eq\_\allowbreak{}joint}: the intermediate facts include \texttt{hw}, \texttt{hpos}, \texttt{hlog}, in source order; the source steps apply or rewrite \texttt{integral\_\allowbreak{}congr\_\allowbreak{}ae}, \texttt{ae\_\allowbreak{}of\_\allowbreak{}all}, \texttt{mul\_\allowbreak{}comm}, \texttt{gradient\_\allowbreak{}bounded\allowbreak{}Log\allowbreak{}Slope\_\allowbreak{}comp}, \texttt{hlog.\allowbreak{}differentiable}. The \texttt{calc} block checks the successive equalities or inequalities in the argument. The remaining scalar rearrangements are checked by \texttt{ring}.
\end{leanbox}


Earlier results used here: \cite[Section~3.26.3]{LangevinLean2026ULASC}.

\section{\texorpdfstring{The bounded time chain rule without a spatial time-majorant}{The bounded time chain rule without a spatial time-majorant}}


On any \(0<a\le b\), the weighted primitive has derivative \(\chi_m(\log r_t)\dot p_t\). The kernel theorem proves \(\dot p\in L^1([a,b]\times E)\), while \(\chi_m\) is bounded by \((m+1)\pi/2\). Thus the derivative of the weighted primitive is jointly integrable.



Integrability at time \(a\) follows from the bounded-entropy lemma after rewriting Gibbs measure as \(J(y)dy\). The generic spacetime FTC/\allowbreak{}Fubini theorem then yields

\[H_m(b)-H_m(a)=\int_a^b q_m(t)\,dt.\]

Applying the same identity on every subinterval expresses \(H_m\) as a constant plus the primitive of the integrable flux \(q_m\). Thus \(H_m\) is absolutely continuous, and Lebesgue differentiation gives \(H_m'(t)=q_m(t)\) almost everywhere. Boundedness of the test and space-time integrability of the density derivative suffice; no common spatial bound for every time derivative is assumed.



Integrate the pointwise chain rule in time first, using the bounded test to dominate the time derivative. Fubini then yields the entropy endpoint identity without requiring an integrable spatial envelope uniform near each individual observation.

\begin{leanbox}{Lean: \texttt{interpolationBoundedEntropy\_\allowbreak{}time\_\allowbreak{}chain\_\allowbreak{}rule}}
\leansource{\texttt{interpolationBoundedEntropy\_\allowbreak{}time\_\allowbreak{}chain\_\allowbreak{}rule} \quad UnadjustedLangevin/\allowbreak{}LSI/\allowbreak{}GaussianEntropyIdentity.lean, lines 301–365}
\begin{leancode}
\leanline{0}{theorem\ interpolationBoundedEntropy\_time\_chain\_rule}
\leanline{4}{(U\ :\ E\ →\ ℝ)\ (μ\ :\ Measure\ E)\ [IsProbabilityMeasure\ μ]}
\leanline{4}{(α\ :\ ℝ)\ (hα\ :\ 0\ <\ α)\ (L\ :\ ℝ≥0)\ (hU\ :\ Differentiable\ ℝ\ U)}
\leanline{4}{(hLip\ :\ LipschitzWith\ L\ (gradient\ U))}
\leanline{4}{(hLSI\ :\ IsLogSobolev\ (ULA.gibbsMeasure\ U)\ α)}
\leanline{4}{(hgrad\ :\ MemLp\ (gradient\ U)\ 2\ μ)\ (a\ b\ :\ ℝ)\ (ha\ :\ 0\ <\ a)\ (hab\ :\ a\ ≤\ b)\ (m\ :\ ℕ)\ :}
\leanline{4}{(interpolationBoundedEntropy\ U\ μ\ m\ b\ -\ interpolationBoundedEntropy\ U\ μ\ m\ a\ =}
\leanline{6}{∫\ t\ in\ a..b,\ interpolationBoundedEntropyFlux\ U\ μ\ m\ t)\ ∧}
\leanline{6}{AbsolutelyContinuousOnInterval\ (interpolationBoundedEntropy\ U\ μ\ m)\ a\ b\ ∧}
\leanline{6}{∀ᵐ\ t\ ∂volume,\ t\ ∈\ Ioo\ a\ b\ →}
\leanline{8}{HasDerivAt\ (interpolationBoundedEntropy\ U\ μ\ m)}
\leanline{10}{(interpolationBoundedEntropyFlux\ U\ μ\ m\ t)\ t}
\end{leancode}
\leanlabel{How the Lean proof follows the argument}
\texttt{interpolation\allowbreak{}Bounded\allowbreak{}Entropy\_\allowbreak{}time\_\allowbreak{}chain\_\allowbreak{}rule}: the intermediate facts include \texttt{h\allowbreak{}G}, \texttt{f}, \texttt{fdot}, \texttt{h\allowbreak{}AC}, \texttt{h\allowbreak{}AE}, in source order; the source steps apply or rewrite \texttt{gibbs\allowbreak{}Measure\_\allowbreak{}eq\_\allowbreak{}with\allowbreak{}Density\_\allowbreak{}real}, \texttt{ENNReal.\allowbreak{}to\allowbreak{}Real\_\allowbreak{}of\allowbreak{}Real}, \texttt{ULA.\allowbreak{}integral\_\allowbreak{}gibbs\_\allowbreak{}eq\_\allowbreak{}weighted}, \texttt{ENNReal.\allowbreak{}to\allowbreak{}Real\_\allowbreak{}inv}, \texttt{integral\_\allowbreak{}const\_\allowbreak{}mul}. The remaining scalar rearrangements are checked by \texttt{linarith}.
\end{leanbox}


Earlier results used here: \cite[Section~3.26.3]{LangevinLean2026ULASC}; \cite[Section~3.46.3]{LangevinLean2026ULASC}.

\section{\texorpdfstring{Expressing the flux on one independent probability space}{Expressing the flux on one independent probability space}}


Take \(X\sim\mu\) and \(Z\sim N(0,I)\) independently, and set \(Y_t=X-tg(X)+\sqrt{2t}Z\). Define

\[V_t(X,Z)=g(X)-\frac{\sqrt{2t}}{2t}Z,\qquad \widehat S_t(X,Z)=S_t(Y_t),\qquad q(t)=-\mathbb E\langle V_t,\widehat S_t\rangle.\]

These fields are defined on the same independent initial-state and Gaussian-noise probability for every time. The frozen velocity is the component spatial flux divided by the positive transition density and evaluated at the proposal. Using this fixed reference space makes joint measurability and time domination available for the limiting entropy identity.



The Gaussian pair-law identity replaces kernel-weighted starting-observation integrals by expectations over one fixed independent state-noise probability. This puts all energy and pairing estimates on the same reference space.

\begin{leanbox}{Lean: \texttt{frozenGaussianVelocityField}, \texttt{interpolationScoreAlongUpdate}, \texttt{interpolationRawEntropyFlux}}
\leansource{\texttt{frozenGaussianVelocityField} \quad UnadjustedLangevin/\allowbreak{}LSI/\allowbreak{}GaussianEntropyIdentity.lean, lines 369–370}
\begin{leancode}
\leanline{0}{def\ frozenGaussianVelocityField\ (U\ :\ E\ →\ ℝ)\ (t\ :\ ℝ)\ (p\ :\ E\ ×\ E)\ :\ E\ :=}
\leanline{2}{gradient\ U\ p.1\ -\ ((1\ /\ (2\ *\ t))\ *\ Real.sqrt\ (2\ *\ t))\ •\ p.2}
\end{leancode}
\leansource{\texttt{interpolationScoreAlongUpdate} \quad UnadjustedLangevin/\allowbreak{}LSI/\allowbreak{}GaussianEntropyIdentity.lean, lines 373–374}
\begin{leancode}
\leanline{0}{def\ interpolationScoreAlongUpdate\ (U\ :\ E\ →\ ℝ)\ (μ\ :\ Measure\ E)\ (t\ :\ ℝ)\ (p\ :\ E\ ×\ E)\ :\ E\ :=}
\leanline{2}{interpolationRelativeScoreField\ U\ μ\ t\ (ULA.malaProposal\ U\ t\ p.1\ p.2)}
\end{leancode}
\leansource{\texttt{interpolationRawEntropyFlux} \quad UnadjustedLangevin/\allowbreak{}LSI/\allowbreak{}GaussianEntropyIdentity.lean, lines 377–379}
\begin{leancode}
\leanline{0}{def\ interpolationRawEntropyFlux\ (U\ :\ E\ →\ ℝ)\ (μ\ :\ Measure\ E)\ (t\ :\ ℝ)\ :\ ℝ\ :=}
\leanline{2}{-(∫\ p\ :\ E\ ×\ E,\ inner\ ℝ\ (frozenGaussianVelocityField\ U\ t\ p)}
\leanline{4}{(interpolationScoreAlongUpdate\ U\ μ\ t\ p)\ ∂μ.prod\ (stdGaussian\ E))}
\end{leancode}
\leanlabel{How the Lean proof follows the argument}
\texttt{frozen\allowbreak{}Gaussian\allowbreak{}Velocity\allowbreak{}Field} is a definition with defining expression \texttt{gradient U p.\allowbreak{}1 -\allowbreak{} ((1 /\allowbreak{} (2 * t)) * Real.\allowbreak{}sqrt (2 * t)) \ensuremath{\bullet} p.\allowbreak{}2}. The right-hand side specifies the object used in the mathematical statement.



\texttt{interpolation\allowbreak{}Score\allowbreak{}Along\allowbreak{}Update} is a definition with defining expression \texttt{interpolation\allowbreak{}Relative\allowbreak{}Score\allowbreak{}Field U \ensuremath{\mu} t (ULA.\allowbreak{}mala\allowbreak{}Proposal U t p.\allowbreak{}1 p.\allowbreak{}2)}. The right-hand side specifies the object used in the mathematical statement.



\texttt{interpolation\allowbreak{}Raw\allowbreak{}Entropy\allowbreak{}Flux} is a definition with defining expression \texttt{-\allowbreak{}(\ensuremath{\int} p : E \ensuremath{\times} E, inner \ensuremath{\mathbb R} (frozen\allowbreak{}Gaussian\allowbreak{}Velocity\allowbreak{}Field U t p)}. The right-hand side specifies the object used in the mathematical statement.
\end{leanbox}


Earlier results used here: \cite[Section~3.34.1]{LangevinLean2026ULASC}.

\section{\texorpdfstring{Square-integrability of the Gaussian flux factors}{Square-integrability of the Gaussian flux factors}}


The initial gradient belongs to \(L^2(\mu)\), and the standard Gaussian noise belongs to \(L^2\). Pulling both to the independent product space and taking their scalar combination proves \(V_t\in L^2(\mu\otimes\gamma)\). The relative score belongs to \(L^2(\rho_t)\). Since \(\rho_t\) is the pushforward of \(\mu\otimes\gamma\) by \(Y_t\), the measurable-map characterization of \(L^2\) gives

\[\widehat S_t=S_t\circ Y_t\in L^2(\mu\otimes\gamma).\]

The initial-gradient and Gaussian moment bounds justify the linear combination defining \(V_t\) in the product \(L^2\) space. The explicit score equals \(\nabla\log r_t\) at positive time, so pushforward change of variables preserves its square energy under composition with \(Y_t\). Cauchy–Schwarz then proves integrability of the inner product of these two fields.



The relative score has finite square energy by the Fisher estimate. The frozen velocity has finite square energy by its Gaussian component and initial gradient, so their flux pairing is absolutely integrable by Cauchy–Schwarz.

\begin{leanbox}{Lean: \texttt{frozenGaussianVelocityField\_\allowbreak{}memLp}, \texttt{interpolationScoreAlongUpdate\_\allowbreak{}memLp}}
\leansource{\texttt{frozenGaussianVelocityField\_\allowbreak{}memLp} \quad UnadjustedLangevin/\allowbreak{}LSI/\allowbreak{}GaussianEntropyIdentity.lean, lines 381–390}
\begin{leancode}
\leanline{0}{theorem\ frozenGaussianVelocityField\_memLp}
\leanline{4}{(U\ :\ E\ →\ ℝ)\ (μ\ :\ Measure\ E)\ [IsProbabilityMeasure\ μ]}
\leanline{4}{(hgrad\ :\ MemLp\ (gradient\ U)\ 2\ μ)\ (t\ :\ ℝ)\ :}
\leanline{4}{MemLp\ (frozenGaussianVelocityField\ U\ t)\ 2\ (μ.prod\ (stdGaussian\ E))}
\end{leancode}
\leansource{\texttt{interpolationScoreAlongUpdate\_\allowbreak{}memLp} \quad UnadjustedLangevin/\allowbreak{}LSI/\allowbreak{}GaussianEntropyIdentity.lean, lines 392–403}
\begin{leancode}
\leanline{0}{theorem\ interpolationScoreAlongUpdate\_memLp}
\leanline{4}{(U\ :\ E\ →\ ℝ)\ (μ\ :\ Measure\ E)\ [IsProbabilityMeasure\ μ]}
\leanline{4}{(L\ :\ ℝ≥0)\ (hU\ :\ Differentiable\ ℝ\ U)\ (hLip\ :\ LipschitzWith\ L\ (gradient\ U))}
\leanline{4}{[IsProbabilityMeasure\ (ULA.gibbsMeasure\ U)]}
\leanline{4}{(hgrad\ :\ MemLp\ (gradient\ U)\ 2\ μ)\ (t\ :\ ℝ)\ (ht\ :\ 0\ <\ t)\ :}
\leanline{4}{MemLp\ (interpolationScoreAlongUpdate\ U\ μ\ t)\ 2\ (μ.prod\ (stdGaussian\ E))}
\end{leancode}
\leanlabel{How the Lean proof follows the argument}
\texttt{frozen\allowbreak{}Gaussian\allowbreak{}Velocity\allowbreak{}Field\_\allowbreak{}mem\allowbreak{}Lp}: the intermediate fact is \texttt{hz}, in source order; the source steps apply or rewrite \texttt{hgrad.\allowbreak{}comp\_\allowbreak{}fst}, \texttt{hz.\allowbreak{}comp\_\allowbreak{}snd}, \texttt{const\_\allowbreak{}smul}, \texttt{Real.\allowbreak{}sqrt}.



\texttt{interpolation\allowbreak{}Score\allowbreak{}Along\allowbreak{}Update\_\allowbreak{}mem\allowbreak{}Lp}: the intermediate fact is \texttt{hs}, in source order; the source steps apply or rewrite \texttt{interpolation\allowbreak{}Relative\allowbreak{}Score\_\allowbreak{}mem\allowbreak{}Lp}, \texttt{interpolation\allowbreak{}Relative\allowbreak{}Score\allowbreak{}Field\_\allowbreak{}eq}, \texttt{mem\allowbreak{}Lp\_\allowbreak{}map\_\allowbreak{}measure\_\allowbreak{}iff}, \texttt{hs.\allowbreak{}aestrongly\allowbreak{}Measurable}.
\end{leanbox}


Earlier results used here: \cite[Section~3.26.3]{LangevinLean2026ULASC}; \cite[Section~3.36.5]{LangevinLean2026ULASC}.

\section{\texorpdfstring{The exact frozen velocity energy}{The exact frozen velocity energy}}


Independence and the zero Gaussian mean cancel the drift-noise cross term:

\[\mathbb E\|V_t\|^2=\mathbb E\|g(X)\|^2+\frac{d}{2t},\qquad t>0.\]

The proof specializes the already verified affine-Gaussian moment formula with coefficients \(1\) and \(-\sqrt{2t}/(2t)\). The squared coefficient is \((2t)/(4t^2)=1/(2t)\), using \((\sqrt{2t})^2=2t\). This exact finite energy gives a compact-time bound that is uniform away from zero. No moment of \(X\) itself is used.



The frozen velocity consists of the Gaussian component score plus the initial drift. Their mixed expectation vanishes by centered independent Gaussian noise, giving the exact sum of component-score and initial-gradient energies.

\begin{leanbox}{Lean: \texttt{frozenGaussianVelocityField\_\allowbreak{}energy}}
\leansource{\texttt{frozenGaussianVelocityField\_\allowbreak{}energy} \quad UnadjustedLangevin/\allowbreak{}LSI/\allowbreak{}GaussianEntropyIdentity.lean, lines 405–418}
\begin{leancode}
\leanline{0}{theorem\ frozenGaussianVelocityField\_energy}
\leanline{4}{(U\ :\ E\ →\ ℝ)\ (μ\ :\ Measure\ E)\ [IsProbabilityMeasure\ μ]}
\leanline{4}{(hgrad\ :\ MemLp\ (gradient\ U)\ 2\ μ)\ (t\ :\ ℝ)\ (ht\ :\ 0\ <\ t)\ :}
\leanline{4}{(∫\ p\ :\ E\ ×\ E,\ ‖frozenGaussianVelocityField\ U\ t\ p‖\ \textasciicircum{}\ 2\ ∂μ.prod\ (stdGaussian\ E))\ =}
\leanline{6}{(∫\ x,\ ‖gradient\ U\ x‖\ \textasciicircum{}\ 2\ ∂μ)\ +\ (Module.finrank\ ℝ\ E\ :\ ℝ)\ /\ (2\ *\ t)}
\end{leancode}
\leanlabel{How the Lean proof follows the argument}
\texttt{frozen\allowbreak{}Gaussian\allowbreak{}Velocity\allowbreak{}Field\_\allowbreak{}energy}: the intermediate fact is \texttt{heq}, in source order; the source steps apply or rewrite \texttt{sub\_\allowbreak{}eq\_\allowbreak{}add\_\allowbreak{}neg}, \texttt{integral\_\allowbreak{}norm\_\allowbreak{}sq\_\allowbreak{}affine\_\allowbreak{}std\allowbreak{}Gaussian}, \texttt{neg\_\allowbreak{}sq}, \texttt{mul\_\allowbreak{}pow}, \texttt{Real.\allowbreak{}sq\_\allowbreak{}sqrt}. The remaining scalar rearrangements are checked by \texttt{field\_\allowbreak{}simp}.
\end{leanbox}

\section{\texorpdfstring{Transporting the bounded flux to independent noise}{Transporting the bounded flux to independent noise}}


The law of \((X,Y_t)\) has density \(K_t(x,y)\) with respect to \(\mu(dx)dy\). Therefore the joint flux identity becomes

\[q_m(t)=-\mathbb E\!\left[w_m(\log r_t(Y_t))\langle V_t,\widehat S_t\rangle\right].\]

The measurable map \((x,z)\mapsto(x,Y_t(x,z))\) has the initial/\allowbreak{}output pair law as its pushforward. This law has density \(K_t\) relative to \(\mu(dx)dy\). The change-of-variables formula for the integrable weighted pairing therefore equates the kernel-weighted integral with the expectation over the independent product space. Finally \(Y_t-(X-tg(X))=\sqrt{2t}Z\) reduces the residual to the specified velocity field. The change of variables concerns the Gaussian law and does not assume an invariant Euler step.



Changing variables in the bounded flux retains the same state-space test evaluated at the proposal. This lets the derivative weight remain bounded by one while the independent-noise energy majorant is applied.

\begin{leanbox}{Lean: \texttt{interpolationBoundedEntropyFlux\_\allowbreak{}eq\_\allowbreak{}raw\_\allowbreak{}weighted}}
\leansource{\texttt{interpolationBoundedEntropyFlux\_\allowbreak{}eq\_\allowbreak{}raw\_\allowbreak{}weighted} \quad UnadjustedLangevin/\allowbreak{}LSI/\allowbreak{}GaussianEntropyIdentity.lean, lines 423–474}
\begin{leancode}
\leanline{0}{theorem\ interpolationBoundedEntropyFlux\_eq\_raw\_weighted}
\leanline{4}{(U\ :\ E\ →\ ℝ)\ (μ\ :\ Measure\ E)\ [IsProbabilityMeasure\ μ]}
\leanline{4}{(L\ :\ ℝ≥0)\ (hU\ :\ Differentiable\ ℝ\ U)\ (hLip\ :\ LipschitzWith\ L\ (gradient\ U))}
\leanline{4}{[IsProbabilityMeasure\ (ULA.gibbsMeasure\ U)]}
\leanline{4}{(hgrad\ :\ MemLp\ (gradient\ U)\ 2\ μ)\ (t\ :\ ℝ)\ (ht\ :\ 0\ <\ t)\ (m\ :\ ℕ)\ :}
\leanline{4}{interpolationBoundedEntropyFlux\ U\ μ\ m\ t\ =}
\leanline{6}{-(∫\ p\ :\ E\ ×\ E,\ boundedLogWeight\ m\ (Real.log\ (interpolationRelativeDensity\ U\ μ\ t}
\leanline{10}{(ULA.malaProposal\ U\ t\ p.1\ p.2)))\ *}
\leanline{8}{inner\ ℝ\ (frozenGaussianVelocityField\ U\ t\ p)}
\leanline{10}{(interpolationScoreAlongUpdate\ U\ μ\ t\ p)\ ∂μ.prod\ (stdGaussian\ E))}
\end{leancode}
\leanlabel{How the Lean proof follows the argument}
\texttt{interpolation\allowbreak{}Bounded\allowbreak{}Entropy\allowbreak{}Flux\_\allowbreak{}eq\_\allowbreak{}raw\_\allowbreak{}weighted}: the intermediate facts include \texttt{T}, \texttt{K}, \texttt{F}, \texttt{h\allowbreak{}N}, \texttt{hchange}, in source order; the source steps apply or rewrite \texttt{frozen\allowbreak{}Gaussian\allowbreak{}Joint\allowbreak{}Law\_\allowbreak{}eq\_\allowbreak{}density}, \texttt{ENNReal.\allowbreak{}to\allowbreak{}Real\_\allowbreak{}of\allowbreak{}Real}, \texttt{ULA.\allowbreak{}mala\allowbreak{}Proposal}, \texttt{add\_\allowbreak{}sub\_\allowbreak{}cancel\_\allowbreak{}left}, \texttt{smul\_\allowbreak{}smul}. The \texttt{calc} block checks the successive equalities or inequalities in the argument. The remaining scalar rearrangements are checked by \texttt{ring}.
\end{leanbox}


Earlier results used here: \cite[Section~3.26.3]{LangevinLean2026ULASC}; \cite[Section~3.34.1]{LangevinLean2026ULASC}; \cite[Section~3.36.5]{LangevinLean2026ULASC}; \cite[Section~3.37.4]{LangevinLean2026ULASC}.

\section{\texorpdfstring{Removing the truncation at fixed positive time}{Removing the truncation at fixed positive time}}


The raw pairing \(\langle V_t,\widehat S_t\rangle\) is integrable by the two \(L^2\) estimates. The scalar map \((x,z)\mapsto\log r_t(Y_t)\) is measurable. Since \(0\le w_m\le1\) and \(w_m\to1\), dominated convergence on the independent probability space yields

\[q_m(t)\longrightarrow q(t),\qquad t>0.\]

The dominating function is \(|\langle V_t,\widehat S_t\rangle|\). This limit is performed after spatial integration by parts; it does not ask for integrability of \(|\log r_t|\,|\dot p_t|\), which is not supplied by the original assumptions.



At a fixed positive time, all factors except the bounded-log weight remain unchanged. Since that weight tends to one and its absolute value is at most one, the integrable full flux dominates the entire sequence.

\begin{leanbox}{Lean: \texttt{interpolationBoundedEntropyFlux\_\allowbreak{}tendsto}}
\leansource{\texttt{interpolationBoundedEntropyFlux\_\allowbreak{}tendsto} \quad UnadjustedLangevin/\allowbreak{}LSI/\allowbreak{}GaussianEntropyIdentity.lean, lines 478–495}
\begin{leancode}
\leanline{0}{theorem\ interpolationBoundedEntropyFlux\_tendsto}
\leanline{4}{(U\ :\ E\ →\ ℝ)\ (μ\ :\ Measure\ E)\ [IsProbabilityMeasure\ μ]}
\leanline{4}{(L\ :\ ℝ≥0)\ (hU\ :\ Differentiable\ ℝ\ U)\ (hLip\ :\ LipschitzWith\ L\ (gradient\ U))}
\leanline{4}{[IsProbabilityMeasure\ (ULA.gibbsMeasure\ U)]}
\leanline{4}{(hgrad\ :\ MemLp\ (gradient\ U)\ 2\ μ)\ (t\ :\ ℝ)\ (ht\ :\ 0\ <\ t)\ :}
\leanline{4}{Tendsto\ (fun\ m\ :\ ℕ\ ↦\ interpolationBoundedEntropyFlux\ U\ μ\ m\ t)\ atTop}
\leanline{6}{({\leanfallback 𝓝}\ (interpolationRawEntropyFlux\ U\ μ\ t))}
\end{leancode}
\leanlabel{How the Lean proof follows the argument}
\texttt{interpolation\allowbreak{}Bounded\allowbreak{}Entropy\allowbreak{}Flux\_\allowbreak{}tendsto}: the intermediate facts include \texttt{hi}, \texttt{hz}, \texttt{hd}, in source order; the source steps apply or rewrite \texttt{interpolation\allowbreak{}Bounded\allowbreak{}Entropy\allowbreak{}Flux\_\allowbreak{}eq\_\allowbreak{}raw\_\allowbreak{}weighted}.
\end{leanbox}


Earlier results used here: \cite[Section~3.26.3]{LangevinLean2026ULASC}; \cite[Section~3.34.1]{LangevinLean2026ULASC}; \cite[Section~3.36.5]{LangevinLean2026ULASC}.

\section{\texorpdfstring{A fixed-reference formula for the relative Fisher energy}{A fixed-reference formula for the relative Fisher energy}}


Define

\[I(t)=\int r_t(y)\|S_t(y)\|^2\,d\pi(y).\]

The relative-density identity gives \(I(t)=\int\|S_t\|^2\,d\rho_t\). Applying the Gaussian pushforward once more yields

\[\mathbb E\|\widehat S_t\|^2=I(t).\]

The first definition fixes the reference measure for time measurability and spacetime integrability. The equality uses the weighted representation of relative Fisher information and the measurable-map Bochner integral formula. The proof supplies strong measurability of the score norm square explicitly, fixing time in the joint measurable score field.



Pushing the relative-score square through the proposal map expresses \(I_t\) as an expectation on the independent-noise source. This representation is jointly measurable and directly comparable with the discrepancy energy.

\begin{leanbox}{Lean: \texttt{interpolationRelativeFisherEnergy}, \texttt{interpolationScoreAlongUpdate\_\allowbreak{}energy}}
\leansource{\texttt{interpolationRelativeFisherEnergy} \quad UnadjustedLangevin/\allowbreak{}LSI/\allowbreak{}GaussianEntropyIdentity.lean, lines 498–500}
\begin{leancode}
\leanline{0}{def\ interpolationRelativeFisherEnergy\ (U\ :\ E\ →\ ℝ)\ (μ\ :\ Measure\ E)\ (t\ :\ ℝ)\ :\ ℝ\ :=}
\leanline{2}{∫\ y,\ interpolationRelativeDensity\ U\ μ\ t\ y\ *}
\leanline{4}{‖interpolationRelativeScoreField\ U\ μ\ t\ y‖\ \textasciicircum{}\ 2\ ∂ULA.gibbsMeasure\ U}
\end{leancode}
\leansource{\texttt{interpolationScoreAlongUpdate\_\allowbreak{}energy} \quad UnadjustedLangevin/\allowbreak{}LSI/\allowbreak{}GaussianEntropyIdentity.lean, lines 502–516}
\begin{leancode}
\leanline{0}{theorem\ interpolationScoreAlongUpdate\_energy}
\leanline{4}{(U\ :\ E\ →\ ℝ)\ (μ\ :\ Measure\ E)\ [IsProbabilityMeasure\ μ]}
\leanline{4}{(L\ :\ ℝ≥0)\ (hU\ :\ Differentiable\ ℝ\ U)\ (hLip\ :\ LipschitzWith\ L\ (gradient\ U))}
\leanline{4}{[IsProbabilityMeasure\ (ULA.gibbsMeasure\ U)]}
\leanline{4}{(hgrad\ :\ MemLp\ (gradient\ U)\ 2\ μ)\ (t\ :\ ℝ)\ (ht\ :\ 0\ <\ t)\ :}
\leanline{4}{(∫\ p\ :\ E\ ×\ E,\ ‖interpolationScoreAlongUpdate\ U\ μ\ t\ p‖\ \textasciicircum{}\ 2\ ∂μ.prod\ (stdGaussian\ E))\ =}
\leanline{6}{interpolationRelativeFisherEnergy\ U\ μ\ t}
\end{leancode}
\leanlabel{How the Lean proof follows the argument}
\texttt{interpolation\allowbreak{}Relative\allowbreak{}Fisher\allowbreak{}Energy} is a definition with defining expression \texttt{\ensuremath{\int} y, interpolation\allowbreak{}Relative\allowbreak{}Density U \ensuremath{\mu} t y *}. The right-hand side specifies the object used in the mathematical statement.



\texttt{interpolation\allowbreak{}Score\allowbreak{}Along\allowbreak{}Update\_\allowbreak{}energy}: the intermediate fact is \texttt{hm}, in source order; the source steps apply or rewrite \texttt{integral\_\allowbreak{}map\_\allowbreak{}of\_\allowbreak{}strongly\allowbreak{}Measurable}.
\end{leanbox}


Earlier results used here: \cite[Section~3.26.3]{LangevinLean2026ULASC}; \cite[Section~3.36.5]{LangevinLean2026ULASC}.

\section{\texorpdfstring{A uniform bound for every truncated flux}{A uniform bound for every truncated flux}}


For any vectors \(v,s\) and any \(0\le w\le1\),

\[|w\langle v,s\rangle|\le\|v\|\|s\|\le\tfrac12(\|v\|^2+\|s\|^2).\]

The second inequality is \((\|v\|-\|s\|)^2\ge0\). Integrate on the independent Gaussian space, use the norm inequality for Bochner integrals, and substitute the two exact energies:

\[|q_m(t)|\le\frac12\left(G_0+\frac d{2t}+I(t)\right),\qquad G_0=\int\|g\|^2\,d\mu.\]

The weighted pairing is integrable because \(w_m\) is bounded and both factors are square integrable. The quadratic upper bound is also integrable by those same two energy estimates. Integral monotonicity therefore justifies the displayed estimate for every truncation before a limiting entropy derivative is known to exist.



The absolute truncated pairing is bounded by half the sum of the two unweighted square energies. This bound contains no truncation index and is therefore suitable for both pointwise-time and time-integrated limits.

\begin{leanbox}{Lean: \texttt{norm\_\allowbreak{}interpolationBoundedEntropyFlux\_\allowbreak{}le}}
\leansource{\texttt{norm\_\allowbreak{}interpolationBoundedEntropyFlux\_\allowbreak{}le} \quad UnadjustedLangevin/\allowbreak{}LSI/\allowbreak{}GaussianEntropyIdentity.lean, lines 520–562}
\begin{leancode}
\leanline{0}{theorem\ norm\_interpolationBoundedEntropyFlux\_le}
\leanline{4}{(U\ :\ E\ →\ ℝ)\ (μ\ :\ Measure\ E)\ [IsProbabilityMeasure\ μ]}
\leanline{4}{(L\ :\ ℝ≥0)\ (hU\ :\ Differentiable\ ℝ\ U)\ (hLip\ :\ LipschitzWith\ L\ (gradient\ U))}
\leanline{4}{[IsProbabilityMeasure\ (ULA.gibbsMeasure\ U)]}
\leanline{4}{(hgrad\ :\ MemLp\ (gradient\ U)\ 2\ μ)\ (t\ :\ ℝ)\ (ht\ :\ 0\ <\ t)\ (m\ :\ ℕ)\ :}
\leanline{4}{‖interpolationBoundedEntropyFlux\ U\ μ\ m\ t‖\ ≤}
\leanline{6}{((∫\ x,\ ‖gradient\ U\ x‖\ \textasciicircum{}\ 2\ ∂μ)\ +\ (Module.finrank\ ℝ\ E\ :\ ℝ)\ /\ (2\ *\ t)\ +}
\leanline{8}{interpolationRelativeFisherEnergy\ U\ μ\ t)\ /\ 2}
\end{leancode}
\leanlabel{How the Lean proof follows the argument}
\texttt{norm\_\allowbreak{}interpolation\allowbreak{}Bounded\allowbreak{}Entropy\allowbreak{}Flux\_\allowbreak{}le}: the intermediate facts include \texttt{P}, \texttt{V}, \texttt{S}, \texttt{hw}, \texttt{hb}, in source order; the source steps apply or rewrite \texttt{Real.\allowbreak{}norm\_\allowbreak{}eq\_\allowbreak{}abs}, \texttt{abs\_\allowbreak{}of\_\allowbreak{}nonneg}, \texttt{interpolation\allowbreak{}Bounded\allowbreak{}Entropy\allowbreak{}Flux\_\allowbreak{}eq\_\allowbreak{}raw\_\allowbreak{}weighted}, \texttt{norm\_\allowbreak{}neg}, \texttt{frozen\allowbreak{}Gaussian\allowbreak{}Velocity\allowbreak{}Field\_\allowbreak{}energy}. The remaining scalar rearrangements are checked by \texttt{nlinarith}.
\end{leanbox}


Earlier results used here: \cite[Section~3.26.3]{LangevinLean2026ULASC}; \cite[Section~3.34.1]{LangevinLean2026ULASC}; \cite[Section~3.36.5]{LangevinLean2026ULASC}.

\section{\texorpdfstring{Joint measurability of both time fluxes}{Joint measurability of both time fluxes}}


The bounded flux is a parameter integral of \(\chi_m(\log r_t(y))\dot p_t(y)\). Both factors are jointly measurable in \((t,y)\), so parameter integration makes \(t\mapsto q_m(t)\) measurable. The raw flux uses the fixed independent space: the map \((t,x,z)\mapsto(t,Y_t)\) is measurable, the explicit relative-score field is jointly strongly measurable, and the velocity field is a measurable algebraic combination of \(g(X)\), \(t^{-1}\), \(\sqrt{2t}\), and \(Z\). Their inner product and its parameter integral are measurable. The definitions at nonpositive times serve only to make these globally specified measurable functions; all analytic conclusions are restricted to positive times.



The chosen score and velocity fields are jointly measurable, and the bounded-log weights are continuous functions of measurable densities. Their products and parameter integrals therefore define measurable scalar fluxes in time.

\begin{leanbox}{Lean: \texttt{interpolationBoundedEntropyFlux\_\allowbreak{}measurable}, \texttt{interpolationRawEntropyFlux\_\allowbreak{}measurable}}
\leansource{\texttt{interpolationBoundedEntropyFlux\_\allowbreak{}measurable} \quad UnadjustedLangevin/\allowbreak{}LSI/\allowbreak{}GaussianEntropyIdentity.lean, lines 564–571}
\begin{leancode}
\leanline{0}{theorem\ interpolationBoundedEntropyFlux\_measurable}
\leanline{4}{(U\ :\ E\ →\ ℝ)\ (μ\ :\ Measure\ E)\ [IsProbabilityMeasure\ μ]}
\leanline{4}{(L\ :\ ℝ≥0)\ (hU\ :\ Differentiable\ ℝ\ U)\ (hLip\ :\ LipschitzWith\ L\ (gradient\ U))\ (m\ :\ ℕ)\ :}
\leanline{4}{Measurable\ (interpolationBoundedEntropyFlux\ U\ μ\ m)}
\end{leancode}
\leansource{\texttt{interpolationRawEntropyFlux\_\allowbreak{}measurable} \quad UnadjustedLangevin/\allowbreak{}LSI/\allowbreak{}GaussianEntropyIdentity.lean, lines 573–586}
\begin{leancode}
\leanline{0}{theorem\ interpolationRawEntropyFlux\_measurable}
\leanline{4}{(U\ :\ E\ →\ ℝ)\ (μ\ :\ Measure\ E)\ [IsProbabilityMeasure\ μ]}
\leanline{4}{(L\ :\ ℝ≥0)\ (hLip\ :\ LipschitzWith\ L\ (gradient\ U))\ :}
\leanline{4}{Measurable\ (interpolationRawEntropyFlux\ U\ μ)}
\end{leancode}
\leanlabel{How the Lean proof follows the argument}
\texttt{interpolation\allowbreak{}Bounded\allowbreak{}Entropy\allowbreak{}Flux\_\allowbreak{}measurable}: the intermediate facts include \texttt{h\ensuremath{\chi}}, \texttt{hh}, in source order; the source steps apply or rewrite \texttt{hh.\allowbreak{}strongly\allowbreak{}Measurable.\allowbreak{}integral\_\allowbreak{}prod\_\allowbreak{}right.\allowbreak{}measurable}.



\texttt{interpolation\allowbreak{}Raw\allowbreak{}Entropy\allowbreak{}Flux\_\allowbreak{}measurable}: the intermediate facts include \texttt{hg}, \texttt{h\allowbreak{}Y}, \texttt{h\allowbreak{}S}, \texttt{h\allowbreak{}V}, in source order; the source steps apply or rewrite \texttt{h\allowbreak{}V.\allowbreak{}inner}, \texttt{integral\_\allowbreak{}prod\_\allowbreak{}right.\allowbreak{}neg.\allowbreak{}measurable}.
\end{leanbox}


Earlier results used here: \cite[Section~3.34.1]{LangevinLean2026ULASC}.

\section{\texorpdfstring{The integrable majorant for the time limit}{The integrable majorant for the time limit}}


Fix \(0<a\le b\). The Gaussian energy estimate proves \(I\in L^1([a,b])\). Set

\[D(t)=\tfrac12\left(G_0+\frac d{2a}+I(t)\right).\]

The interval has finite Lebesgue measure, so the constant term is integrable; hence \(D\in L^1([a,b])\). Since \(t\ge a>0\), \(d/(2t)\le d/(2a)\). The uniform flux estimate gives \(|q_m(t)|\le D(t)\) for every \(m\). The fixed-time limit \(q_m(t)\to q(t)\), followed by continuity of the absolute value, also gives \(|q(t)|\le D(t)\). Thus one integrable time function bounds all truncated fluxes and their pointwise limit, as required for dominated convergence in time.



On a compact positive interval, the score and velocity square energies have a common integrable time bound. Their half-sum dominates every truncated flux and the limiting flux, establishing the spacetime dominated-convergence hypothesis.

\begin{leanbox}{Lean: \texttt{interpolation\_\allowbreak{}entropy\_\allowbreak{}flux\_\allowbreak{}time\_\allowbreak{}domination}}
\leansource{\texttt{interpolation\_\allowbreak{}entropy\_\allowbreak{}flux\_\allowbreak{}time\_\allowbreak{}domination} \quad UnadjustedLangevin/\allowbreak{}LSI/\allowbreak{}GaussianEntropyIdentity.lean, lines 590–620}
\begin{leancode}
\leanline{0}{theorem\ interpolation\_entropy\_flux\_time\_domination}
\leanline{4}{(U\ :\ E\ →\ ℝ)\ (μ\ :\ Measure\ E)\ [IsProbabilityMeasure\ μ]}
\leanline{4}{(L\ :\ ℝ≥0)\ (hU\ :\ Differentiable\ ℝ\ U)\ (hLip\ :\ LipschitzWith\ L\ (gradient\ U))}
\leanline{4}{[IsProbabilityMeasure\ (ULA.gibbsMeasure\ U)]}
\leanline{4}{(hgrad\ :\ MemLp\ (gradient\ U)\ 2\ μ)\ (a\ b\ :\ ℝ)\ (ha\ :\ 0\ <\ a)\ :}
\leanline{4}{∃\ D\ :\ ℝ\ →\ ℝ,\ Integrable\ D\ (volume.restrict\ (Ioc\ a\ b))\ ∧}
\leanline{6}{(∀\ m,\ ∀ᵐ\ t\ ∂volume.restrict\ (Ioc\ a\ b),}
\leanline{8}{‖interpolationBoundedEntropyFlux\ U\ μ\ m\ t‖\ ≤\ D\ t)\ ∧}
\leanline{6}{∀ᵐ\ t\ ∂volume.restrict\ (Ioc\ a\ b),\ ‖interpolationRawEntropyFlux\ U\ μ\ t‖\ ≤\ D\ t}
\end{leancode}
\leanlabel{How the Lean proof follows the argument}
\texttt{interpolation\_\allowbreak{}entropy\_\allowbreak{}flux\_\allowbreak{}time\_\allowbreak{}domination}: the intermediate facts include \texttt{G}, \texttt{D}, \texttt{h\allowbreak{}G}, \texttt{htpos}, \texttt{hlim}, in source order; the source steps apply or rewrite \texttt{norm\_\allowbreak{}interpolation\allowbreak{}Bounded\allowbreak{}Entropy\allowbreak{}Flux\_\allowbreak{}le}, \texttt{ht.\allowbreak{}1.\allowbreak{}le}, \texttt{le\_\allowbreak{}of\_\allowbreak{}tendsto}, \texttt{Eventually.\allowbreak{}of\_\allowbreak{}forall}.
\end{leanbox}


Earlier results used here: \cite[Section~3.26.3]{LangevinLean2026ULASC}.

\section{\texorpdfstring{The integrated entropy identity}{The integrated entropy identity}}


Let \(H(t)=H(\rho_t\mid\pi)\), represented by the finite relative entropy. For each truncation,

\[H_m(b)-H_m(a)=\int_a^b q_m(t)\,dt.\]

At both endpoints, the spatial dominated-convergence theorem gives \(H_m\to H\). In time, the common integrable majorant \(D\) and the pointwise convergence \(q_m\to q\) give \(\int_a^b q_m\to\int_a^b q\). Uniqueness of limits therefore yields

\[H(b)-H(a)=\int_a^b q(t)\,dt.\]

For the time integral, use Lebesgue measure restricted to \((a,b]\); endpoint singletons have measure zero. The endpoint limits were proved under Gibbs measure, and the flux limit uses the integrable time bound \(D\). These distinct dominated-convergence arguments give the identity without assuming temporal regularity of the full entropy.



Take the truncation limit in the two endpoint values and in the integrated flux separately. Their established dominations then yield an exact identity for the relative entropy between any two positive times.

\begin{leanbox}{Lean: \texttt{actualInterpolationEntropy}, \texttt{interpolationEntropy\_\allowbreak{}integral\_\allowbreak{}identity}}
\leansource{\texttt{actualInterpolationEntropy} \quad UnadjustedLangevin/\allowbreak{}LSI/\allowbreak{}GaussianEntropyIdentity.lean, lines 623–624}
\begin{leancode}
\leanline{0}{def\ actualInterpolationEntropy\ (U\ :\ E\ →\ ℝ)\ (μ\ :\ Measure\ E)\ (t\ :\ ℝ)\ :\ ℝ\ :=}
\leanline{2}{(InformationTheory.klDiv\ (interpolationLaw\ U\ μ\ t)\ (ULA.gibbsMeasure\ U)).toReal}
\end{leancode}
\leansource{\texttt{interpolationEntropy\_\allowbreak{}integral\_\allowbreak{}identity} \quad UnadjustedLangevin/\allowbreak{}LSI/\allowbreak{}GaussianEntropyIdentity.lean, lines 628–655}
\begin{leancode}
\leanline{0}{theorem\ interpolationEntropy\_integral\_identity}
\leanline{4}{(U\ :\ E\ →\ ℝ)\ (μ\ :\ Measure\ E)\ [IsProbabilityMeasure\ μ]}
\leanline{4}{(α\ :\ ℝ)\ (hα\ :\ 0\ <\ α)\ (L\ :\ ℝ≥0)\ (hU\ :\ Differentiable\ ℝ\ U)}
\leanline{4}{(hLip\ :\ LipschitzWith\ L\ (gradient\ U))}
\leanline{4}{(hLSI\ :\ IsLogSobolev\ (ULA.gibbsMeasure\ U)\ α)}
\leanline{4}{(hgrad\ :\ MemLp\ (gradient\ U)\ 2\ μ)\ (a\ b\ :\ ℝ)\ (ha\ :\ 0\ <\ a)\ (hab\ :\ a\ ≤\ b)\ :}
\leanline{4}{actualInterpolationEntropy\ U\ μ\ b\ -\ actualInterpolationEntropy\ U\ μ\ a\ =}
\leanline{6}{∫\ t\ in\ a..b,\ interpolationRawEntropyFlux\ U\ μ\ t}
\end{leancode}
\leanlabel{How the Lean proof follows the argument}
\texttt{actual\allowbreak{}Interpolation\allowbreak{}Entropy} is a definition with defining expression \texttt{(Information\allowbreak{}Theory.\allowbreak{}kl\allowbreak{}Div (interpolation\allowbreak{}Law U \ensuremath{\mu} t) (ULA.\allowbreak{}gibbs\allowbreak{}Measure U)).\allowbreak{}to\allowbreak{}Real}. The right-hand side specifies the object used in the mathematical statement.



\texttt{interpolation\allowbreak{}Entropy\_\allowbreak{}integral\_\allowbreak{}identity}: the intermediate fact is \texttt{h}, in source order; the source steps apply or rewrite \texttt{interval\allowbreak{}Integral.\allowbreak{}integral\_\allowbreak{}of\_\allowbreak{}le}, \texttt{integral\_\allowbreak{}identity\_\allowbreak{}of\_\allowbreak{}dominated\_\allowbreak{}limits}, \texttt{interpolation\allowbreak{}Bounded\allowbreak{}Entropy\allowbreak{}Flux\_\allowbreak{}tendsto}, \texttt{ha.\allowbreak{}trans}, \texttt{ht.\allowbreak{}1}.
\end{leanbox}


Earlier results used here: \cite[Section~3.26.3]{LangevinLean2026ULASC}.

\section{\texorpdfstring{Absolute continuity and the almost-everywhere entropy derivative}{Absolute continuity and the almost-everywhere entropy derivative}}


The raw flux is measurable and satisfies \(|q|\le D\), so it is integrable on every positive-time compact interval. Apply the full entropy identity with terminal time \(t\in[a,b]\):

\[H(t)=H(a)+\int_a^t q(u)\,du.\]

An indefinite integral of an \(L^1\) function is absolutely continuous, and the Lebesgue differentiation theorem identifies its derivative almost everywhere. Thus

\[H\in AC([a,b]),\qquad H'(t)=q(t)\quad\text{for almost every }t\in(a,b).\]

The proof derives interval integrability from the raw-flux measurability and the verified majorant, then invokes absolute continuity and the almost-everywhere derivative of an integrable time primitive. This proves the required time regularity from the stated potential, log-Sobolev, and finite initial score-energy assumptions. A later score-pairing identity rewrites the raw Gaussian flux as negative Fisher energy plus the frozen gradient discrepancy, yielding the quantitative entropy-production estimate.



The flux is integrable on each positive compact interval, so the exact integrated identity makes entropy locally absolutely continuous there. Its derivative equals that flux almost everywhere, supplying the precise temporal regularity used in the final scalar proof.

\begin{leanbox}{Lean: \texttt{interpolationEntropy\_\allowbreak{}local\_\allowbreak{}chain\_\allowbreak{}rule}}
\leansource{\texttt{interpolationEntropy\_\allowbreak{}local\_\allowbreak{}chain\_\allowbreak{}rule} \quad UnadjustedLangevin/\allowbreak{}LSI/\allowbreak{}GaussianEntropyIdentity.lean, lines 660–677}
\begin{leancode}
\leanline{0}{theorem\ interpolationEntropy\_local\_chain\_rule}
\leanline{4}{(U\ :\ E\ →\ ℝ)\ (μ\ :\ Measure\ E)\ [IsProbabilityMeasure\ μ]}
\leanline{4}{(α\ :\ ℝ)\ (hα\ :\ 0\ <\ α)\ (L\ :\ ℝ≥0)\ (hU\ :\ Differentiable\ ℝ\ U)}
\leanline{4}{(hLip\ :\ LipschitzWith\ L\ (gradient\ U))}
\leanline{4}{(hLSI\ :\ IsLogSobolev\ (ULA.gibbsMeasure\ U)\ α)}
\leanline{4}{(hgrad\ :\ MemLp\ (gradient\ U)\ 2\ μ)\ (a\ b\ :\ ℝ)\ (ha\ :\ 0\ <\ a)\ (hab\ :\ a\ ≤\ b)\ :}
\leanline{4}{AbsolutelyContinuousOnInterval\ (actualInterpolationEntropy\ U\ μ)\ a\ b\ ∧}
\leanline{6}{∀ᵐ\ t\ ∂volume,\ t\ ∈\ Ioo\ a\ b\ →}
\leanline{8}{HasDerivAt\ (actualInterpolationEntropy\ U\ μ)\ (interpolationRawEntropyFlux\ U\ μ\ t)\ t}
\end{leancode}
\leanlabel{How the Lean proof follows the argument}
\texttt{bounded\allowbreak{}Entropy\_\allowbreak{}tendsto} proves the endpoint limit of the bounded logarithmic primitives. \texttt{interpolation\allowbreak{}Entropy\_\allowbreak{}local\_\allowbreak{}chain\_\allowbreak{}rule} obtains \texttt{hq}, integrability of the raw flux from \texttt{interpolation\_\allowbreak{}entropy\_\allowbreak{}flux\_\allowbreak{}time\_\allowbreak{}domination}, and applies \texttt{regularity\_\allowbreak{}of\_\allowbreak{}integral\_\allowbreak{}identity} to the integrated identity. The result contains local absolute continuity and an almost-everywhere \texttt{Has\allowbreak{}Deriv\allowbreak{}At} statement. \texttt{interpolation\allowbreak{}Entropy\_\allowbreak{}derivative\_\allowbreak{}bound\_\allowbreak{}ae} applies this on a countable family of intervals starting at reciprocal integers. \texttt{hall} gathers their almost-everywhere statements, and \texttt{exists\_\allowbreak{}nat\_\allowbreak{}one\_\allowbreak{}div\_\allowbreak{}lt} chooses an interval containing any given positive time. Rewriting by \texttt{hd.\allowbreak{}deriv} and applying \texttt{interpolation\allowbreak{}Raw\allowbreak{}Entropy\allowbreak{}Flux\_\allowbreak{}le} gives the stated production inequality almost everywhere.
\end{leanbox}


Earlier results used here: \cite[Section~3.26.3]{LangevinLean2026ULASC}.

\begin{thebibliography}{ULA-LSI}
\bibitem[ULA-SC]{LangevinLean2026ULASC}
LangevinLean Project. \emph{Unadjusted Langevin under Strong Convexity}. LL-01, 2026. Draft report, content revision d50c0f64815b. \url{https://odakyildiz.github.io/LangevinLean/docs/pdf/UnadjustedLangevin/FullProof.pdf}.
\end{thebibliography}

\end{document}
